\begin{verso}{Jeremías 9:23–24}
Así dijo el \SENOR: No se alabe el sabio en su sabiduría, ni en su valentía se alabe el valiente, ni el rico se alabe en sus riquezas. Mas alábese en esto el que se hubiere de alabar: en entenderme y conocerme, que yo soy el \SENOR, que hago misericordia, juicio y justicia en la tierra; porque estas cosas quiero, dice el \SENOR.
\end{verso}

\section{Una gloria fuera de lugar}

Este libro nace de un versículo. Jeremías lo pronunció ante un pueblo que estaba a punto de perderlo todo. Judá confiaba en la sabiduría de sus consejeros, en la fuerza de sus alianzas y en la prosperidad de sus ciudades, mientras el juicio se acercaba. En ese contexto, Dios no desarma a su pueblo con un argumento, sino con una pregunta implícita: ¿en qué te estás gloriando?

El profeta nombra las tres grandes seguridades humanas: el saber, el poder y el tener. El sabio confía en lo que entiende; el valiente, en lo que puede; el rico, en lo que posee. Ninguna de ellas es mala en sí misma, pero ninguna puede sostener el peso de la vida. Por eso Dios no prohíbe gloriarse; cambia el \emph{objeto} de la gloria. Hay una sola cosa en la que el ser humano puede gloriarse sin engañarse: entender y conocer a Dios.

Y el Dios que debe ser conocido no es una idea abstracta. Él mismo se describe por lo que hace y por lo que ama: misericordia, juicio y justicia en la tierra. En el versículo que abre este libro ya aparecen, entonces, los atributos de Dios. Conocer a Dios es conocerle en sus perfecciones, tal como Él las ha revelado.

Arthur W. Pink abrió su conocida obra sobre los atributos de Dios con este mismo texto, y afirmó que “un conocimiento espiritual y salvador de Dios es la mayor necesidad de toda criatura humana” (Pink, 2020, p. 4). Inmediatamente después escribió una frase que justifica todo lo que sigue: “No se puede confiar, ni servir, ni adorar a un Dios desconocido” (Pink, 2020, p. 4). Si la fe, el servicio y la adoración dependen de conocer a Dios, estudiar sus atributos no es un lujo para especialistas: es una necesidad para todo creyente.

\section{Entender y conocer}

Jeremías une dos verbos que no significan lo mismo. \destacar{Entender} a Dios es comprender cómo actúa en la historia del ser humano y en la vida de cada persona, qué valora según su carácter santo y fiel, y cómo gobierna con justicia, misericordia y verdad. Es un conocimiento que pasa por la mente y que lleva a discernir su voluntad y a vivir conforme a ella. \destacar{Conocer} a Dios, en cambio, es vivir en relación personal con Él: una comunión de pacto, real y cotidiana, que se expresa en fidelidad y en una vida delante de su rostro.

Los dos son necesarios. Una teología sin relación se vuelve fría y distante; una experiencia sin entendimiento se vuelve frágil y confusa. Jeremías no pide escoger entre la mente y el corazón: pide que ambos se rindan ante el mismo Dios.

La historia de la medicina ofrece una ilustración de esta diferencia. En 1984, el médico australiano Barry J. Marshall quería demostrar que una bacteria, \emph{Helicobacter pylori}, causaba la inflamación del estómago. Sus colegas no aceptaban la hipótesis, así que Marshall la probó en su propio cuerpo: ingirió la bacteria y desarrolló gastritis (Marshall et al., 1985). Aquella experiencia confirmó su teoría y cambió la medicina digestiva; años después recibió el Premio Nobel de Medicina (2005). Marshall no solo \emph{sabía} de la bacteria: la conoció en carne propia. Así es el conocimiento de Dios que pide la Escritura. Pink lo dijo con claridad: “Necesitamos algo más que un conocimiento teórico de Dios” (Pink, 2020, p. 4).

El profeta Oseas expresa la misma idea como una decisión que se renueva:

\begin{verso}{Oseas 6:3}
Y conoceremos, y proseguiremos en conocer al \SENOR; como el alba está dispuesta su salida, y vendrá a nosotros como la lluvia, como la lluvia tardía y temprana a la tierra.
\end{verso}

Conocer a Dios no se agota en un momento. Es una búsqueda que prosigue, con la misma certeza con que cada mañana llega el alba. El Nuevo Testamento lleva esta idea a su punto más alto, porque identifica el conocimiento de Dios con la vida misma:

\begin{verso}{Juan 17:3}
Y esta es la vida eterna: que te conozcan a ti, el único Dios verdadero, y a Jesucristo, a quien has enviado.
\end{verso}

\section{Un Dios infinito ante una mente limitada}

Antes de avanzar es necesario reconocer un límite. Dios es infinito; nosotros, criaturas finitas. Nuestra mente puede recibir lo que Él revela, pero no puede abarcarlo por completo. Ningún estudio, por cuidadoso que sea, agotará lo que Dios es. Por eso este libro se escribe desde la \destacar{humildad}: lo que aquí se afirma de Dios es verdadero en la medida en que Él lo ha revelado, pero siempre será parcial.

Isaías lo expresa con una imagen que no admite discusión:

\begin{verso}{Isaías 55:8–9}
Porque mis pensamientos no son vuestros pensamientos, ni vuestros caminos mis caminos, dijo el \SENOR. Como son más altos los cielos que la tierra, así son mis caminos más altos que vuestros caminos, y mis pensamientos más que vuestros pensamientos.
\end{verso}

Si la distancia es tan grande, ¿cómo podemos conocer algo de Dios? El Nuevo Testamento responde que el conocimiento de Dios no nace de nuestra capacidad, sino de su Espíritu. Solo Dios se conoce a sí mismo de manera perfecta, y es su Espíritu quien nos da a conocer lo que de otro modo nos sería inaccesible. Este es el versículo que ancla todo el estudio:

\begin{verso}{1 Corintios 2:10–11}
Pero Dios nos las reveló a nosotros por el Espíritu; porque el Espíritu todo lo escudriña, aun lo profundo de Dios. Porque ¿quién de los hombres sabe las cosas del hombre, sino el espíritu del hombre que está en él? Así tampoco nadie conoció las cosas de Dios, sino el Espíritu de Dios.
\end{verso}

Pablo establece aquí dos verdades a la vez. La primera es un límite: “nadie conoció las cosas de Dios, sino el Espíritu de Dios”. Ninguna mente humana puede escudriñar a Dios desde fuera. La segunda es un don: “Dios nos las reveló a nosotros por el Espíritu”. Lo que no podemos alcanzar por nosotros mismos, Dios lo comunica. Estudiar los atributos de Dios es, entonces, recibir con gratitud lo que el Espíritu ha revelado, sin pretender ir más allá.

Pink, al cerrar su obra, insiste en esta misma actitud: “la incomprensibilidad de la naturaleza divina debería enseñarnos humildad, precaución y reverencia” (Pink, 2020, p. 100). Y cita al puritano John Howe, quien comparó nuestro conocimiento de Dios con “tener un breve resumen de un libro muy voluminoso o […] observar un pequeño paisaje de un país gigantesco. Él nos ha dado un verdadero informe de Sí mismo, pero no completo” (Pink, 2020, p. 100). Incluso en la eternidad seguirá siendo así; Pink recoge las palabras de John Dick: “Para comprender la perfección infinita, deben volverse infinitos ellos mismos” (Pink, 2020, p. 101).

De esta convicción se desprenden dos aclaraciones sobre la naturaleza de este libro.

La primera es que este libro es una \destacar{recopilación de estudios individuales}. Cada capítulo nació como un estudio particular sobre un atributo de Dios, preparado y expuesto por separado. Aquí se reúnen y se ordenan, pero conservan su origen: son estudios que buscan conocer a Dios desde un aspecto concreto de su revelación.

La segunda es que \destacar{no se abarcarán todos los atributos de Dios}. Un Dios infinito tiene perfecciones que exceden cualquier lista. Este libro estudia algunas de las que la Escritura revela con mayor claridad, sin pretender que sean las únicas. Habrá siempre más en Dios de lo que podemos nombrar, y quizá perfecciones suyas que nunca llegaremos a comprender. Lejos de desalentarnos, esta realidad nos invita a adorar: el Dios que estudiamos es más grande que nuestro estudio.

\precision{la clasificación y la selección de atributos de este libro son herramientas para ordenar lo revelado; no pretenden describir a Dios de manera exhaustiva.}

\section{La perspectiva de este libro}

Este libro estudia los atributos de Dios desde tres convicciones.

La primera es una \destacar{mirada hebrea}. El Dios de la Biblia no se presentó a Israel con una lista de definiciones, sino por su Nombre, por sus actos en la historia y por su pacto. Israel conoció su fidelidad al salir de Egipto, su santidad en el Sinaí y su misericordia en el desierto. Por eso cada atributo se estudiará a partir de los textos y de los encuentros en los que Dios se dio a conocer a su pueblo.

La segunda es una \destacar{cosmovisión paulina}. Pablo fue un judío formado en las Escrituras de Israel que, al encontrarse con Cristo, las reconoció cumplidas en Él. Él mismo lo recordó ante sus compatriotas:

\begin{verso}{Hechos 22:3}
Yo de cierto soy judío, nacido en Tarso de Cilicia, pero criado en esta ciudad, instruido a los pies de Gamaliel, estrictamente conforme a la ley de nuestros padres, celoso de Dios, como hoy lo sois todos vosotros.
\end{verso}

Pablo no abandonó las Escrituras de Israel para hablar de Dios; las leyó a la luz de Cristo. Cuando contempla la obra de Dios en la historia de la salvación, cita a Isaías y a Job, y termina en adoración:

\begin{verso}{Romanos 11:33, 36}
¡Oh profundidad de las riquezas de la sabiduría y de la ciencia de Dios! ¡Cuán insondables son sus juicios, e inescrutables sus caminos! […] Porque de él, y por él, y para él, son todas las cosas. A él sea la gloria por los siglos. Amén.
\end{verso}

Esa es la actitud que este libro quiere imitar: estudiar con rigor, pero terminar adorando.

La tercera convicción es un método \destacar{expositivo}. Cada atributo se desarrollará a partir de textos bíblicos concretos, explicados en su contexto, y no desde definiciones abstractas. La obra de A. W. Pink acompañará el recorrido como apoyo, pero el punto de partida será siempre la Escritura.

\precision{la unión de una mirada hebrea con una cosmovisión paulina es el criterio de lectura propio de este libro; no pretende ser la única lectura posible de los textos.}

\section{Dios se da a conocer}

El Dios de la Escritura no puede ser descubierto por la sola investigación humana; se da a conocer. Pink lo afirmó sin rodeos: “Tal Dios no puede ser descubierto mediante investigación. Sólo puede ser conocido cuando es revelado al corazón por el Espíritu Santo, a través de la Palabra” (Pink, 2020, p. 8).

Esa revelación llega por varios caminos. La \destacar{naturaleza} muestra su existencia, su poder, su gloria y su sabiduría; es una revelación general, al alcance de todos (Sal 19:1–6; Ro 1:19–20). La vida misma, con sus \destacar{criaturas y experiencias}, muestra su cuidado, su bondad y su gobierno continuo (Job 12:7–10; Sal 104). Pero es la \destacar{Palabra} la que revela con claridad su carácter, su voluntad y su plan de redención (Sal 19:7–11).

\begin{verso}{Salmo 19:1}
Los cielos cuentan la gloria de Dios, y el firmamento anuncia la obra de sus manos.
\end{verso}

\begin{verso}{Romanos 1:20}
Porque las cosas invisibles de él, su eterno poder y deidad, se hacen claramente visibles desde la creación del mundo, siendo entendidas por medio de las cosas hechas, de modo que no tienen excusa.
\end{verso}

La creación deja al ser humano sin excusa, porque muestra que Dios existe y que es poderoso. Sin embargo, para conocer su santidad, su misericordia, su fidelidad y su plan de salvación, necesitamos que Él hable. Y lo hizo de manera plena en Jesucristo, como se verá más adelante.

\section{El Nombre: la colección de sus atributos}

El pasaje que mejor explica qué entendemos por atributo está en el Sinaí. Israel acababa de fabricar el becerro de oro, y Moisés, intercediendo por el pueblo, hizo la petición más audaz de la Escritura:

\begin{verso}{Éxodo 33:18–19}
Él entonces dijo: Te ruego que me muestres tu gloria. Y le respondió: Yo haré pasar todo mi bien delante de tu rostro, y proclamaré el nombre del \SENOR{} delante de ti; y tendré misericordia del que tendré misericordia, y seré clemente para con el que seré clemente.
\end{verso}

Moisés pidió ver la gloria, y Dios respondió proclamando su Nombre. Pink recoge esta conexión en una frase que se ha vuelto el principio de este libro: “El nombre es la colección de sus atributos” (Pink, 2020, p. 100). La proclamación llega en el capítulo siguiente:

\begin{verso}{Éxodo 34:6–7}
Y pasando el \SENOR{} por delante de él, proclamó: ¡\SENOR! ¡\SENOR! fuerte, misericordioso y piadoso; tardo para la ira, y grande en misericordia y verdad; que guarda misericordia a millares, que perdona la iniquidad, la rebelión y el pecado, y que de ningún modo tendrá por inocente al malvado; que visita la iniquidad de los padres sobre los hijos y sobre los hijos de los hijos, hasta la tercera y cuarta generación.
\end{verso}

Dios no se definió con categorías abstractas. Se dio a conocer proclamando su Nombre, y en ese Nombre aparecen juntos su poder, su misericordia, su gracia, su paciencia, su amor leal, su verdad y su justicia. La misericordia no anula la justicia, ni la justicia apaga la misericordia: ambas forman parte del mismo Nombre. El propio Antiguo Testamento vuelve una y otra vez a esta proclamación (Nm 14:18; Neh 9:17; Sal 86:15; 103:8; Jl 2:13; Jon 4:2), señal de que Israel la recibió como la descripción central de su Dios.

\section{¿Qué es un atributo?}

Un \destacar{atributo} es una perfección eterna que Dios revela de sí mismo en la Escritura y por la cual podemos conocerle. No es una parte de Dios, como si Él estuviera compuesto de piezas, sino Dios mismo considerado desde uno de los aspectos de su ser. Dios no \emph{tiene} amor y santidad como quien posee objetos; Dios \emph{es} amor y \emph{es} santo, y lo es completamente y a la vez (Berkhof, 1938).

De esta definición se desprenden tres principios que guiarán todo el libro.

El primero es que \destacar{todos los atributos son eternos}. Dios no adquiere perfecciones con el tiempo ni las pierde. Lo que Él es hoy, lo ha sido siempre y lo será siempre.

\begin{verso}{Malaquías 3:6}
Porque yo el \SENOR{} no cambio; por esto, hijos de Jacob, no habéis sido consumidos.
\end{verso}

El segundo es que \destacar{los atributos no se oponen entre sí}. La justicia de Dios no compite con su misericordia, ni su santidad con su amor. En Dios actúan en perfecta armonía, como lo canta el salmista:

\begin{verso}{Salmo 85:10}
La misericordia y la verdad se encontraron; la justicia y la paz se besaron.
\end{verso}

El tercero es que \destacar{los atributos no se agotan}. Lo que conocemos de Dios es verdadero, pero no exhaustivo. Job, después de describir el poder de Dios sobre la creación, reconoce que apenas ha tocado la orilla:

\begin{verso}{Job 26:14}
He aquí, estas cosas son sólo los bordes de sus caminos; ¡y cuán leve es el susurro que hemos oído de él! Pero el trueno de su poder, ¿quién lo puede comprender?
\end{verso}

Estudiar los atributos de Dios exige, por tanto, dos actitudes a la vez: confianza, porque lo que Él ha revelado es verdadero; y humildad, porque lo que conocemos es apenas un susurro de lo que Él es.

\section{Primera clasificación: Ser, Identidad y Relación}

Toda clasificación de los atributos es una herramienta para ordenar el estudio; la Escritura no presenta una lista sistemática. Este libro los organiza en tres grupos, a partir de tres preguntas: ¿qué es Dios?, ¿quién es Dios? y ¿cómo se revela Dios en su relación con el ser humano?

\subsection{Atributos propios de su Ser}

Son las perfecciones que describen lo que Dios \emph{es} en su ser interior, eterno e independiente de toda la creación. No dependen de nada externo: Dios era así antes de crear el mundo y lo seguiría siendo aunque nada hubiera sido creado. La Escritura los expresa con la fórmula más directa posible: “Dios es…”. Este libro reconoce tres: Dios es \emph{uno}, Dios es \emph{Espíritu} y Dios es \emph{amor}.

\begin{verso}{Deuteronomio 6:4}
Oye, Israel: el \SENOR{} nuestro Dios, el \SENOR{} uno es.
\end{verso}

\begin{verso}{Juan 4:24}
Dios es Espíritu; y los que le adoran, en espíritu y en verdad es necesario que adoren.
\end{verso}

\begin{verso}{1 Juan 4:8}
El que no ama, no ha conocido a Dios; porque Dios es amor.
\end{verso}

\subsection{Atributos propios de su Identidad}

Son las perfecciones que expresan cómo es Dios en su esencia eterna: su carácter propio y absoluto, aquello que lo distingue como el único Dios. Aquí se encuentran su santidad, su eternidad, su soberanía, su justicia, su inmutabilidad, su omnipotencia, su omnipresencia y su omnisciencia. En el centro de todas ellas está la santidad, el atributo que los serafines proclaman tres veces delante de su trono:

\begin{verso}{Isaías 6:3}
Y el uno al otro daba voces, diciendo: Santo, santo, santo, \SENOR{} de los ejércitos; toda la tierra está llena de su gloria.
\end{verso}

\subsection{Atributos de relación con el ser humano}

Aquí se requiere precisión. \destacar{Todos los atributos son perfecciones eternas de Dios}; ninguno comenzó a existir con la creación. Sin embargo, algunos de ellos se fueron \emph{revelando} en la historia de la relación de Dios con el ser humano. La sabiduría se hizo visible al crear el mundo (Pr 3:19); la misericordia y la ira, cuando apareció el pecado (Gn 3; Ro 1:18); la fidelidad, la verdad y el celo, en el pacto con Israel (Éx 20:5; 34:6). Son la sabiduría, la verdad, la misericordia, la fidelidad, el celo y la ira santa, a los que se suman la bondad, la paciencia y la gracia, que Pink estudia en capítulos propios.

\begin{verso}{Salmo 103:8}
Misericordioso y clemente es el \SENOR; lento para la ira, y grande en misericordia.
\end{verso}

\precision{decir que la ira o la misericordia son atributos “de relación” no significa que sean temporales. Son la santidad y el amor eternos de Dios manifestados ante el pecado y la necesidad humana. Lo histórico es su manifestación, no su existencia.}

La Tabla~\ref{tab:mapa} resume esta primera clasificación.

\begin{table}[htbp]
\centering\small
\caption{Primera clasificación de los atributos de Dios}\label{tab:mapa}
\renewcommand{\arraystretch}{1.3}
\begin{tabular}{@{}P{3.4cm} P{3.0cm} P{6.2cm} P{3.0cm}@{}}
\toprule
\textbf{Grupo} & \textbf{Pregunta} & \textbf{Atributos} & \textbf{Textos} \\ \midrule
Propios de su Ser & ¿Qué es? & Uno, Espíritu, Amor & Dt 6:4; Jn 4:24; 1 Jn 4:8 \\
Propios de su Identidad & ¿Quién es? & Santo, Eterno, Soberano, Justo, Inmutable, Omnipotente, Omnipresente, Omnisciente & Is 6:3; Sal 90:2; Sal 115:3; Mal 3:6 \\
De relación con el ser humano & ¿Cómo se revela? & Sabiduría, Verdad, Misericordia, Fidelidad, Celo, Ira santa, Bondad, Paciencia, Gracia & Éx 34:6–7; Sal 103:8–13 \\ \bottomrule
\end{tabular}
\end{table}

\section{Segunda clasificación: comunicables e incomunicables}

La tradición reformada ofrece una segunda forma de ordenar los atributos, complementaria de la anterior: según si Dios los comparte o no con el ser humano (Berkhof, 1938; Grudem, 2007).

Los \destacar{atributos incomunicables}, o no transferibles, son aquellos que solo Dios posee. Ninguna criatura es eterna, inmutable u omnipresente. En ellos Dios es absolutamente único:

\begin{verso}{Isaías 46:9}
Acordaos de las cosas pasadas desde los tiempos antiguos; porque yo soy Dios, y no hay otro Dios, y nada hay semejante a mí…
\end{verso}

Los \destacar{atributos comunicables}, o transferibles, son aquellos que Dios posee en perfección y que refleja, en forma limitada, en el ser humano. El fundamento de esta comunicación es la creación del hombre a imagen de Dios:

\begin{verso}{Génesis 1:27}
Y creó Dios al hombre a su imagen, a imagen de Dios lo creó; varón y hembra los creó.
\end{verso}

Lo comunicable no es solo un privilegio; también es un llamado. Pedro, citando el libro de Levítico, lo convierte en mandato para la vida diaria:

\begin{verso}{1 Pedro 1:15–16}
sino, como aquel que os llamó es santo, sed también vosotros santos en toda vuestra manera de vivir; porque escrito está: Sed santos, porque yo soy santo.
\end{verso}

\begin{table}[htbp]
\centering\small
\caption{Segunda clasificación de los atributos de Dios}\label{tab:cruce}
\renewcommand{\arraystretch}{1.3}
\begin{tabular}{@{}P{4.0cm} P{11.6cm}@{}}
\toprule
\textbf{Incomunicables} & Uno, Eterno, Inmutable, Omnipresente, Omnipotente, Omnisciente \\
\textbf{Comunicables} & Amor, Espíritu, Santo, Justo, Soberano (en forma reflejada), Sabiduría, Verdad, Misericordia, Fidelidad, Celo, Ira santa, Bondad, Paciencia, Gracia \\ \bottomrule
\end{tabular}
\par\smallskip{\footnotesize\color{grisaz}\textit{Nota.} El poder y el conocimiento infinitos no se comunican, pero sí su reflejo finito. Siguiendo a Berkhof (1938), la espiritualidad se cuenta entre los comunicables: el ser humano recibe espíritu (Gn 2:7), no la espiritualidad infinita de Dios.}
\end{table}

\section{La plenitud en Cristo}

Cada capítulo de este libro terminará con un apartado sobre \destacar{cómo se reflejó ese atributo en Cristo}. La razón es sencilla: si el Nombre de Dios es la colección de sus atributos, Cristo es ese Nombre hecho carne. En Él los atributos de Dios no aparecen por partes, sino en plenitud. Pablo lo afirma dos veces en la carta a los Colosenses:

\begin{verso}{Colosenses 1:19–20}
por cuanto agradó al Padre que en él habitase toda plenitud, y por medio de él reconciliar consigo todas las cosas, así las que están en la tierra como las que están en los cielos, haciendo la paz mediante la sangre de su cruz.
\end{verso}

\begin{verso}{Colosenses 2:9}
Porque en él habita corporalmente toda la plenitud de la Deidad,
\end{verso}

La carta a los Hebreos usa otra imagen: Cristo es el resplandor de la gloria de Dios y la imagen misma de su ser.

\begin{verso}{Hebreos 1:3}
el cual, siendo el resplandor de su gloria, y la imagen misma de su sustancia, y quien sustenta todas las cosas con la palabra de su poder, habiendo efectuado la purificación de nuestros pecados por medio de sí mismo, se sentó a la diestra de la Majestad en las alturas,
\end{verso}

Lo que Moisés pidió ver en el Sinaí, los discípulos lo vieron en Jesús. Juan describe al Verbo encarnado con palabras que recuerdan la proclamación de Éxodo 34:6: “lleno de gracia y de verdad”.

\begin{verso}{Juan 1:14}
Y aquel Verbo fue hecho carne, y habitó entre nosotros (y vimos su gloria, gloria como del unigénito del Padre), lleno de gracia y de verdad.
\end{verso}

Por eso Jesús pudo responder a Felipe, cuando este le pidió ver al Padre:

\begin{verso}{Juan 14:9}
Jesús le dijo: ¿Tanto tiempo hace que estoy con vosotros, y no me has conocido, Felipe? El que me ha visto a mí, ha visto al Padre; ¿cómo, pues, dices tú: Muéstranos el Padre?
\end{verso}

Y Pablo une en un solo versículo los grandes temas de esta introducción —la luz, el conocimiento y la gloria— y los sitúa en el rostro de Cristo:

\begin{verso}{2 Corintios 4:6}
Porque Dios, que mandó que de las tinieblas resplandeciese la luz, es el que resplandeció en nuestros corazones, para iluminación del conocimiento de la gloria de Dios en la faz de Jesucristo.
\end{verso}

\section{Nuestra gloria}

La introducción termina donde comenzó, porque el propio Nuevo Testamento cierra el círculo. Toda la Escritura invita a trasladar nuestra gloria desde nosotros mismos hacia Dios: “En Dios nos gloriaremos todo el tiempo, y para siempre alabaremos tu nombre” (Sal 44:8). Pablo retoma esa invitación y la lleva a su cumplimiento en Cristo:

\begin{verso}{Romanos 5:11}
Y no sólo esto, sino que también nos gloriamos en Dios por el Señor nuestro Jesucristo, por quien hemos recibido ahora la reconciliación.
\end{verso}

Más aún, cita las palabras mismas de Jeremías para explicar lo que Dios hizo en Cristo:

\begin{verso}{1 Corintios 1:30–31}
Mas por él estáis vosotros en Cristo Jesús, el cual nos ha sido hecho por Dios sabiduría, justificación, santificación y redención; para que, como está escrito: El que se gloría, gloríese en el Señor.
\end{verso}

La sabiduría, la fuerza y la riqueza en las que el ser humano quería gloriarse, Dios las ha dado en Cristo. Por eso conocer los atributos de Dios no es un ejercicio intelectual: es aprender a gloriarnos en el lugar correcto.

\section{Cuando Dios se revela}

La Escritura muestra que, cuando Dios se revela, no solo habla: transforma. Moisés era un pastor en el desierto cuando vio una zarza que ardía sin consumirse. Allí Dios se le reveló y le dio a conocer su Nombre:

\begin{verso}{Éxodo 3:14}
Y respondió Dios a Moisés: YO SOY EL QUE SOY. Y dijo: Así dirás a los hijos de Israel: YO SOY me envió a vosotros.
\end{verso}

Ese encuentro cambió la vida de Moisés y, por medio de él, la historia de todo un pueblo. El mismo efecto busca este libro. No pretende solo informar sobre Dios, sino conducir a conocerle, porque el conocimiento verdadero de Dios transforma la vida.

\section{Plan de la obra}

El libro reúne los estudios realizados sobre los atributos de Dios y los organiza en tres partes que siguen la primera clasificación, y una conclusión. Como se dijo, no pretende abarcar todos los atributos: estudia algunas de las perfecciones que Dios ha revelado. Cada capítulo desarrollará un atributo con la misma estructura: el texto bíblico que lo ancla, su definición, su enseñanza en diálogo con Pink, un encuentro en la historia de Israel donde ese atributo se reveló, \destacar{su reflejo en Cristo}, y su aplicación a la vida.

\begin{table}[htbp]
\centering\small
\caption{Estructura del libro}\label{tab:plan}
\renewcommand{\arraystretch}{1.3}
\begin{tabular}{@{}P{3.4cm} P{8.6cm} P{3.6cm}@{}}
\toprule
\textbf{Parte} & \textbf{Capítulos} & \textbf{Pink (2020)} \\ \midrule
Introducción & Conocer a Dios en sus atributos & Prefacio; cap. 1 \\
I. Ser & Uno · Espíritu · Amor & Caps. 1, 16–17 \\
II. Identidad & Santo · Eterno · Soberano · Justo · Inmutable · Omnipotente · Omnipresente · Omnisciente & Caps. 2–9 \\
III. Relación & Sabiduría · Verdad · Misericordia · Fidelidad · Celo · Ira santa · Bondad · Paciencia · Gracia & Caps. 10–15, 18 \\
Conclusión & La gloria de Dios: la comunión de sus atributos en Cristo & Cap. 19 \\ \bottomrule
\end{tabular}
\end{table}

\section{Síntesis}

Jeremías enseña que la única gloria legítima del ser humano es entender y conocer a Dios. Ese conocimiento no es una conquista humana, sino un don: Dios se da a conocer proclamando su Nombre, y su Nombre es la colección de sus atributos. Todos ellos son perfecciones eternas; unos expresan lo que Dios es en su Ser, otros quién es en su Identidad, y otros se fueron revelando en su relación con el ser humano. Unos son exclusivamente suyos; otros los refleja en nosotros, creados a su imagen.

Pero la meta de este libro no es una lista. La \destacar{gloria de Dios es la comunión perfecta de todos sus atributos}, que actúan como uno solo porque Dios es uno. Esa plenitud no quedó oculta: habita corporalmente en Cristo, en cuyo rostro brilla el conocimiento de la gloria de Dios. Por eso cada capítulo terminará en Él, y por eso todo el libro es una invitación a cumplir el mandato de Jeremías: \emph{el que se gloría, gloríese en el Señor}.
